\documentclass[10pt]{extarticle}
\usepackage[a4paper,margin=1.5cm]{geometry}
\usepackage[T1]{fontenc}
\usepackage{amsmath,amssymb}
\usepackage{graphicx}
\usepackage[export]{adjustbox}
\usepackage{multicol}
\usepackage{enumitem}
\usepackage{booktabs,array,tabularx}
\setlength{\columnsep}{0.7cm}
\setlength{\parindent}{0pt}
\setlength{\parskip}{3pt}
\pagestyle{empty}
\makeatletter
\newcommand{\qfigdef}[4]{\@namedef{qf@file@#1}{#2}\@namedef{qf@w@#1}{#3}\@namedef{qf@wide@#1}{#4}}
\newcommand{\qf@check}[1]{\@ifundefined{qf@file@#1}{\PackageError{score2}{Unknown figure id #1}{}}{}}
\newcommand{\qf@img}[4][]{\includegraphics[width=#2\dimexpr\@nameuse{qf@w@#4}pt\relax,max width=#3,max height=0.45\textheight,#1]{\@nameuse{qf@file@#4}}}
\newif\ifqf@full
\newif\ifqf@boxed
\newsavebox{\qf@box}
\newcommand{\qf@open}{\global\qf@fulltrue\par\global\qf@boxedtrue
  \setbox\qf@box\vbox\bgroup\hsize\textwidth\linewidth\textwidth}
\newcommand{\qf@emit}[1]{\par\egroup\global\qf@boxedfalse
  \ifdim\dimexpr\ht\qf@box+\dp\qf@box\relax<#1\textheight\noindent\usebox{\qf@box}\par\else\unvbox\qf@box\fi}
\newcommand{\qfullsplit}{\ifqf@boxed\qf@emit{0.85}\fi}
\newcommand{\qf@close}{\ifqf@boxed\qf@emit{0.45}\fi\global\qf@fullfalse}
\newcommand{\qfullbegin}{\end{multicols}\qf@open}
\newcommand{\qfullend}{\qf@close\begin{multicols}{2}}
\newcommand{\qfullstart}{\qf@open}
\newcommand{\qfulldone}{\qf@close}
\newcommand{\qfig}[2][1]{\qf@check{#2}%
  \ifqf@full\qf@blk{#1}{#2}\else
  \ifnum\@nameuse{qf@wide@#2}=1 \expandafter\@firstoftwo\else\expandafter\@secondoftwo\fi
  {\end{multicols}\begin{center}\qf@img{#1}{\textwidth}{#2}\end{center}\begin{multicols}{2}}%
  {\qf@blk{#1}{#2}}\fi}
\newcommand{\qf@blk}[2]{\par\smallskip{\centering\qf@img{#1}{\linewidth}{#2}\par}\smallskip}
\newcommand{\qfigi}[2][1]{\qf@check{#2}\qf@img[valign=c]{#1}{\linewidth}{#2}}
\makeatother
\newcommand{\papersec}[1]{\par\bigskip{\centering\Large\bfseries PAPER-#1\par}\medskip}
\newcommand{\partsec}[2]{\par\medskip\noindent{\large\bfseries PART-#1: #2}\par\smallskip}
\newcommand{\sect}[1]{\par\smallskip\noindent\textbf{#1}\par}
\newcommand{\qn}[1]{\par\smallskip\noindent\textbf{#1.}~}
\newcommand{\paraq}[2]{\par\smallskip\noindent\textbf{Paragraph for Questions #1 and #2}\par}
\newenvironment{qoptions}{\begin{itemize}[leftmargin=1.8em,itemsep=1pt,topsep=2pt,label={}]}{\end{itemize}}
\newcommand{\opt}[2]{\item[(#1)] #2}
% figure registry (generated by score2_tools.py finalize)
\qfigdef{p1_phy_s1_q01_1}{Score-II - Test 30 (27 Mar 2024)-img/p1_phy_s1_q01_1.png}{71.7}{0}
\qfigdef{p1_phy_s1_q07_1}{Score-II - Test 30 (27 Mar 2024)-img/p1_phy_s1_q07_1.png}{53.5}{0}
\qfigdef{p1_phy_s1_q08_2}{Score-II - Test 30 (27 Mar 2024)-img/p1_phy_s1_q08_2.png}{55.8}{0}
\qfigdef{p1_phy_s1_q08_1}{Score-II - Test 30 (27 Mar 2024)-img/p1_phy_s1_q08_1.png}{47.5}{0}
\qfigdef{p1_phy_s1_q08_3}{Score-II - Test 30 (27 Mar 2024)-img/p1_phy_s1_q08_3.png}{48.3}{0}
\qfigdef{p1_phy_s1_q08_4}{Score-II - Test 30 (27 Mar 2024)-img/p1_phy_s1_q08_4.png}{77.2}{0}
\qfigdef{p1_phy_s1_q09_1}{Score-II - Test 30 (27 Mar 2024)-img/p1_phy_s1_q09_1.png}{108.3}{0}
\qfigdef{p1_phy_s1_q11_1}{Score-II - Test 30 (27 Mar 2024)-img/p1_phy_s1_q11_1.png}{95.8}{0}
\qfigdef{p1_phy_s2_q01_1}{Score-II - Test 30 (27 Mar 2024)-img/p1_phy_s2_q01_1.png}{70.4}{0}
\qfigdef{p1_phy_s2_q02_1}{Score-II - Test 30 (27 Mar 2024)-img/p1_phy_s2_q02_1.png}{80.4}{0}
\qfigdef{p1_chem_s1_q06_1}{Score-II - Test 30 (27 Mar 2024)-img/p1_chem_s1_q06_1.png}{76.5}{0}
\qfigdef{p1_chem_s1_q13_1}{Score-II - Test 30 (27 Mar 2024)-img/p1_chem_s1_q13_1.png}{45.8}{0}
\qfigdef{p1_chem_s1_q13_2}{Score-II - Test 30 (27 Mar 2024)-img/p1_chem_s1_q13_2.png}{46.8}{0}
\qfigdef{p1_chem_s1_q13_3}{Score-II - Test 30 (27 Mar 2024)-img/p1_chem_s1_q13_3.png}{47.7}{0}
\qfigdef{p1_chem_s1_q13_4}{Score-II - Test 30 (27 Mar 2024)-img/p1_chem_s1_q13_4.png}{49.9}{0}
\qfigdef{p1_chem_s1_q14_1}{Score-II - Test 30 (27 Mar 2024)-img/p1_chem_s1_q14_1.png}{23.0}{0}
\qfigdef{p1_chem_s1_q14_2}{Score-II - Test 30 (27 Mar 2024)-img/p1_chem_s1_q14_2.png}{63.0}{0}
\qfigdef{p1_chem_s1_q14_3}{Score-II - Test 30 (27 Mar 2024)-img/p1_chem_s1_q14_3.png}{23.0}{0}
\qfigdef{p1_chem_s1_q14_4}{Score-II - Test 30 (27 Mar 2024)-img/p1_chem_s1_q14_4.png}{23.0}{0}
\qfigdef{p1_chem_s1_q15_1}{Score-II - Test 30 (27 Mar 2024)-img/p1_chem_s1_q15_1.png}{145.0}{0}
\qfigdef{p1_chem_s1_q15_2}{Score-II - Test 30 (27 Mar 2024)-img/p1_chem_s1_q15_2.png}{44.9}{0}
\qfigdef{p1_chem_s1_q15_3}{Score-II - Test 30 (27 Mar 2024)-img/p1_chem_s1_q15_3.png}{44.9}{0}
\qfigdef{p1_chem_s1_q16_1}{Score-II - Test 30 (27 Mar 2024)-img/p1_chem_s1_q16_1.png}{137.0}{0}
\qfigdef{p1_chem_s1_q16_2}{Score-II - Test 30 (27 Mar 2024)-img/p1_chem_s1_q16_2.png}{39.9}{0}
\qfigdef{p1_chem_s1_q16_3}{Score-II - Test 30 (27 Mar 2024)-img/p1_chem_s1_q16_3.png}{39.9}{0}
\qfigdef{p1_chem_s1_q16_4}{Score-II - Test 30 (27 Mar 2024)-img/p1_chem_s1_q16_4.png}{37.5}{0}
\qfigdef{p1_chem_s1_q16_5}{Score-II - Test 30 (27 Mar 2024)-img/p1_chem_s1_q16_5.png}{38.0}{0}
\qfigdef{p1_chem_s1_q17_1}{Score-II - Test 30 (27 Mar 2024)-img/p1_chem_s1_q17_1.png}{151.4}{0}
\qfigdef{p1_chem_s1_q20_1}{Score-II - Test 30 (27 Mar 2024)-img/p1_chem_s1_q20_1.png}{101.4}{0}
\qfigdef{p1_chem_s1_q20_2}{Score-II - Test 30 (27 Mar 2024)-img/p1_chem_s1_q20_2.png}{133.0}{0}
\qfigdef{p1_chem_s1_q20_3}{Score-II - Test 30 (27 Mar 2024)-img/p1_chem_s1_q20_3.png}{84.8}{0}
\qfigdef{p1_chem_s1_q20_4}{Score-II - Test 30 (27 Mar 2024)-img/p1_chem_s1_q20_4.png}{94.4}{0}
\qfigdef{p1_chem_s1_q20_5}{Score-II - Test 30 (27 Mar 2024)-img/p1_chem_s1_q20_5.png}{91.1}{0}
\qfigdef{p1_chem_s2_q07_1}{Score-II - Test 30 (27 Mar 2024)-img/p1_chem_s2_q07_1.png}{92.8}{0}
\qfigdef{p1_chem_s2_q09_1}{Score-II - Test 30 (27 Mar 2024)-img/p1_chem_s2_q09_1.png}{151.6}{0}
\qfigdef{p1_chem_s2_q10_1}{Score-II - Test 30 (27 Mar 2024)-img/p1_chem_s2_q10_1.png}{161.1}{0}
\makeatletter\@beginparpenalty=10000\makeatother
\begin{document}
\begin{center}{\Large\bfseries JEE (Main+Advanced) Score-II --- Test 30 (27 Mar 2024)}\end{center}
\begin{multicols}{2}
\raggedcolumns
\partsec{1}{PHYSICS}
\sect{Section-I}

\qn{1} A galvanometer (G) of 2$\Omega$ resistance is connected in the given circuit. The ratio of charge stored in C$_1$ and C$_2$ is :
\qfig{p1_phy_s1_q01_1}
\begin{qoptions}
\opt{A}{$\dfrac23$}
\opt{B}{$\dfrac32$}
\opt{C}{1}
\opt{D}{$\dfrac12$}
\end{qoptions}

\qn{2} A uniform magnetic field of $2\times10^{-3}$ T acts along positive Y-direction. A rectangular loop of sides 20 cm and 10 cm with current of 5 A is Y-Z plane. The current is in anticlockwise sense with reference to negative X axis. Magnitude and direction of the torque is :
\begin{qoptions}
\opt{A}{$2\times10^{-4}$ N $-$ m along positive Z $-$direction}
\opt{B}{$2\times10^{-4}$ N $-$ m along negative Z-direction}
\opt{C}{$2\times10^{-4}$ N $-$ m along positive X-direction}
\opt{D}{$2\times10^{-4}$ N $-$ m along positive Y-direction}
\end{qoptions}

\qn{3} Match List I with List II
\par\nobreak
\begin{center}\small
\begin{tabular}{|c|p{3.3cm}|c|p{1.9cm}|}\hline
 & \multicolumn{1}{c|}{\textbf{List I}} & & \multicolumn{1}{c|}{\textbf{List II}}\\\hline
A. & $\displaystyle\oint\vec B\cdot\vec{dl}=\mu_0i_c+\mu_0\varepsilon_0\dfrac{d\phi_E}{dt}$ & I. & Gauss' law for electricity\\\hline
B. & $\displaystyle\oint\vec E\cdot\vec{dl}=\dfrac{d\phi_B}{dt}$ & II. & Gauss' law for magnetism\\\hline
C. & $\displaystyle\oint\vec E\cdot\vec{dA}=\dfrac{Q}{\varepsilon_0}$ & III. & Faraday law\\\hline
D. & $\displaystyle\oint\vec B\cdot\vec{dA}=0$ & IV. & Ampere $-$ Maxwell law\\\hline
\end{tabular}
\end{center}
Chose the correct answer from the options given below
\begin{qoptions}
\opt{A}{A-IV, B-I, C-III, D-II}
\opt{B}{A-II, B-III, C-I, D-IV}
\opt{C}{A-IV, B-III, C-I, D-II}
\opt{D}{A-I, B-II, C-III, D-IV}
\end{qoptions}

\qn{4} The deflection in moving coil galvanometer falls from 25 divisions to 5 division when a shunt of 24$\Omega$ is applied. The resistance of galvanometer coil will be :
\begin{qoptions}
\opt{A}{12$\Omega$}
\opt{B}{96$\Omega$}
\opt{C}{48$\Omega$}
\opt{D}{100$\Omega$}
\end{qoptions}

\qn{5} A capacitor of capacitance 100$\mu$F is charged to a potential of 12V and connected to a 6.4mH inductor to produce oscillations. The maximum current in the circuit would be :
\begin{qoptions}
\opt{A}{3.2 A}
\opt{B}{1.5 A}
\opt{C}{2.0 A}
\opt{D}{1.2 A}
\end{qoptions}

\qn{6} Primary side of a transformer is connected to 230 V, 50 Hz supply. Turns ratio of primary to secondary winding is 10 : 1. Load resistance connected to secondary side is 46 $\Omega$. The power consumed in it is :
\begin{qoptions}
\opt{A}{12.5 W}
\opt{B}{10.0 W}
\opt{C}{11.5 W}
\opt{D}{12.0 W}
\end{qoptions}

\qn{7} The magnetic field at the centre of a wire loop formed by two semicircular wires of radii\\
$R_1=2\pi$ m and $R_2=4\pi$ m carrying current $I=4$A as per figure given below is $\alpha\times10^{-7}$ T. The value of $\alpha$ is \underline{\hspace{2em}}. (Centre O is common for all segments)
\qfig{p1_phy_s1_q07_1}
\begin{qoptions}
\opt{A}{3}
\opt{B}{4}
\opt{C}{8}
\opt{D}{5}
\end{qoptions}

\qn{8} Which of the following circuits is reverse - biased ?
\begin{qoptions}
\opt{A}{\qfigi{p1_phy_s1_q08_2}}
\opt{B}{\qfigi{p1_phy_s1_q08_1}}
\opt{C}{\qfigi{p1_phy_s1_q08_3}}
\opt{D}{\qfigi{p1_phy_s1_q08_4}}
\end{qoptions}

\qn{9} The truth table of the given circuit diagram is:
\qfig{p1_phy_s1_q09_1}
\begin{qoptions}
\opt{A}{\begin{tabular}[t]{|c|c|c|}\hline A&B&Y\\\hline 0&0&1\\0&1&0\\1&0&0\\1&1&1\\\hline\end{tabular}}
\opt{B}{\begin{tabular}[t]{|c|c|c|}\hline A&B&Y\\\hline 0&0&0\\0&1&1\\1&0&1\\1&1&0\\\hline\end{tabular}}
\opt{C}{\begin{tabular}[t]{|c|c|c|}\hline A&B&Y\\\hline 0&0&0\\0&1&0\\1&0&0\\1&1&1\\\hline\end{tabular}}
\opt{D}{\begin{tabular}[t]{|c|c|c|}\hline A&B&Y\\\hline 0&0&1\\0&1&1\\1&0&1\\1&1&0\\\hline\end{tabular}}
\end{qoptions}

\qn{10} In a photoelectric effect experiment a light of frequency 1.5 times the threshold frequency is made to fall on the surface of photosensitive material. Now if the frequency is halved and intensity is doubled, the number of photo electrons emitted will be:
\begin{qoptions}
\opt{A}{Doubled}
\opt{B}{Quadrupled}
\opt{C}{Zero}
\opt{D}{Halved}
\end{qoptions}

\qn{11} A disc of radius $R$ and mass $M$ is rolling horizontally without slipping with speed $v$. It then moves up an inclined smooth surface as shown in figure. The maximum height that the disc can go up the incline is :
\qfig{p1_phy_s1_q11_1}
\begin{qoptions}
\opt{A}{$\dfrac{v^2}{g}$}
\opt{B}{$\dfrac34\dfrac{v^2}{g}$}
\opt{C}{$\dfrac12\dfrac{v^2}{g}$}
\opt{D}{$\dfrac23\dfrac{v^2}{g}$}
\end{qoptions}

\qn{12} A big drop is formed by coalescing 1000 small droplets of water. The surface energy will become :
\begin{qoptions}
\opt{A}{100 times}
\opt{B}{10 times}
\opt{C}{$\dfrac1{100}$th}
\opt{D}{$\dfrac1{10}$th}
\end{qoptions}

\qn{13} A cricket player catches a ball of mass 120 g moving with 25 m/s speed. If the catching process is completed in 0.1 s then the magnitude of force exerted by the ball on the hand of player will be (in SI unit):
\begin{qoptions}
\opt{A}{24}
\opt{B}{12}
\opt{C}{25}
\opt{D}{30}
\end{qoptions}

\qn{14} If $R$ is the radius of the earth and the acceleration due to gravity on the surface of earth is $g=\pi^2\ \mathrm{m/s^2}$, then the length of the second's pendulum at a height $h=2R$ from the surface of the earth will be :
\begin{qoptions}
\opt{A}{$\dfrac29$ m}
\opt{B}{$\dfrac19$ m}
\opt{C}{$\dfrac49$ m}
\opt{D}{$\dfrac89$ m}
\end{qoptions}

\qn{15} A ball of mass 0.5 kg is attached to a string of length 50 cm. The ball is rotated on a horizontal circular path about its vertical axis. The maximum tension that the string can bear is 400 N. The maximum possible value of angular velocity of the ball in rad/s is,:
\begin{qoptions}
\opt{A}{1600}
\opt{B}{40}
\opt{C}{1000}
\opt{D}{20}
\end{qoptions}

\qn{16} The radius ($r$), length ($\ell$) and resistance ($R$) of a metal wire was measured in the laboratory as\\
$r=(0.35\pm0.05)$ cm\\
$R=(100\pm10)$ ohm\\
$\ell=(15\pm0.2)$ cm\\
The percentage error in resistivity of the material of the wire is :
\begin{qoptions}
\opt{A}{25.6\%}
\opt{B}{39.9\%}
\opt{C}{37.3\%}
\opt{D}{35.6\%}
\end{qoptions}

\qn{17} Identify the physical quantity that cannot be measured using spherometer :
\begin{qoptions}
\opt{A}{Radius of curvature of concave surface}
\opt{B}{Specific rotation of liquids}
\opt{C}{Thickness of thin plates}
\opt{D}{Radius of curvature of convex surface}
\end{qoptions}

\qn{18} The distance between object and its 3 times magnified virtual image as produced by a convex lens is 20 cm. The focal length of the lens used is \underline{\hspace{2em}} cm.
\begin{qoptions}
\opt{A}{15}
\opt{B}{18}
\opt{C}{23}
\opt{D}{25}
\end{qoptions}

\qn{19} The number of air molecules per cm$^3$ increased from $3\times10^{19}$ to $12\times10^{19}$. The ratio of collision frequency of air molecules before and after the increase in number respectively is
\begin{qoptions}
\opt{A}{1.25}
\opt{B}{0.25}
\opt{C}{0.75}
\opt{D}{0.50}
\end{qoptions}

\qn{20} Consider two containers A and B containing monoatomic gases at the same Pressure ($P$), Volume ($V$) and Temperature ($T$). The gas in A is compressed isothermally to $\dfrac18$ of its original volume while the gas B is compressed adiabatically to $\dfrac18$ of its original volume. The ratio of final pressure of gas in B to that of gas in A is :
\begin{qoptions}
\opt{A}{8}
\opt{B}{$8^{\frac32}$}
\opt{C}{$\dfrac18$}
\opt{D}{4}
\end{qoptions}

\sect{Section-II}

\qn{1} In the following circuit, the battery has an emf of 2 V and an internal resistance of $\dfrac23\ \Omega$. The power consumption in the entire circuit is \underline{\hspace{2em}} W.
\qfig{p1_phy_s2_q01_1}

\qn{2} The distance between charges $+q$ and $-q$ is $2\,\ell$ and between $+2$ q and $-2$ q is $4\,\ell$. The electrostatic potential at point P at a distance $r$ from centre O is $-\alpha\left[\dfrac{q\ell}{r^2}\right]\times10^9$V , where the value of $\alpha$ is \underline{\hspace{2em}}.\\
(Use $\dfrac1{4\pi\varepsilon_0}=9\times10^9\ \mathrm{Nm^2C^{-2}}$)
\qfig{p1_phy_s2_q02_1}

\qn{3} In a nuclear fission process, a high mass nuclide (A $\approx$ 236) with binding energy 7.6 MeV/Nucleon dissociated into middle mass nuclides (A $\approx$ 118), having binding energy of 8.6 MeV/Nucleon. The energy released in the process would be \underline{\hspace{2em}} MeV.

\qn{4} A particular hydrogen - like ion emits the radiation of frequency $3\times10^{15}$ Hz when it makes transition from $n=2$ to $n=1$. The frequency of radiation emitted in transition from $n=3$ to $n=1$ is $\dfrac x9\times10^{15}$ Hz,\\
when $x=$ \underline{\hspace{2em}}.

\qn{5} A mass $m$ is suspended from a spring of negligible mass and the system oscillates with a frequency $f_1$. The frequency of oscillations if a mass 9 $m$ is suspended from the same spring is $f_2$. The value of $\dfrac{f_1}{f_2}$ is \underline{\hspace{2em}}.

\qn{6} A particle initially at rest starts moving from reference point. $x=0$ along x-axis, with velocity $v$ that varies as $v=4\sqrt x$ m/s. The acceleration of the particle is \underline{\hspace{2em}} ms$^{-2}$.

\qn{7} A monochromatic light of wavelength 6000\AA\ is incident on the single slit of width 0.01 mm. If the diffraction pattern is formed at the focus of the convex lens of focal length 20 cm, the linear width of the central maximum (in mm) is :

\qn{8} A tuning fork resonates with a sonometer wire of length 1 m stretched with a tension of 6 N. When the tension in the wire is changed to 54 N, the same tuning fork produces 12 beats per second with it. The frequency of the tuning fork is \underline{\hspace{2em}} Hz.

\qn{9} Two plates A and B have thermal conductivities 84 Wm$^{-1}$K$^{-1}$ and 126 Wm$^{-1}$K$^{-1}$ respectively. They have same surface area and same thickness. They are placed in contact along their surfaces. If the temperatures of the outer surfaces of A and B are kept at 100 $^\circ$C and 0 $^\circ$C respectively, then the temperature of the surface of contact in steady state is \underline{\hspace{2em}}$^\circ$C.

\qn{10} A steel rod of length 1 m and cross sectional area $10^{-4}$ m$^2$ is heated from 0$^\circ$C to 200$^\circ$C without being allowed to extend or bend. The compressive tension produced in the rod is \underline{\hspace{2em}} $\times10^4$ N. (Given Young's modulus of steel $=2\times10^{11}\ \mathrm{Nm^{-2}}$, coefficient of linear expansion $=10^{-5}\mathrm{K^{-1}}$)

\partsec{2}{CHEMISTRY}
\sect{Section-I}

\qn{1} Choose the correct option for free expansion of an ideal gas under adiabatic condition from the following :
\begin{qoptions}
\opt{A}{$q=0,\ \Delta T\neq0,\ w=0$}
\opt{B}{$q=0,\ \Delta T<0,\ w\neq0$}
\opt{C}{$q\neq0,\ \Delta T=0,\ w=0$}
\opt{D}{$q=0,\ \Delta T=0,\ w=0$}
\end{qoptions}

\qn{2} Solubility of calcium phosphate (molecular mass, $M$) in water is $W_g$ per 100 mL at 25$^\circ$ C. Its solubility product at 25$^\circ$C will be approximately.
\begin{qoptions}
\opt{A}{$10^7\left(\dfrac WM\right)^3$}
\opt{B}{$10^7\left(\dfrac WM\right)^5$}
\opt{C}{$10^3\left(\dfrac WM\right)^5$}
\opt{D}{$10^5\left(\dfrac WM\right)^5$}
\end{qoptions}

\qn{3} The half-life of radioisotopic bromine - 82 is 36 hours. The fraction which remains after one day is \underline{\hspace{2em}} $\times10^{-2}$.\\
(Given antilog 0.2006 = 1.587)
\begin{qoptions}
\opt{A}{63}
\opt{B}{126}
\opt{C}{32}
\opt{D}{100}
\end{qoptions}

\qn{4} The following concentrations were observed at 500 K for the formation of NH$_3$ from N$_2$ and H$_2$. At equilibrium : [N$_2$] = $2\times10^{-2}$ M,\\
\relax[H$_2$] = $3\times10^{-2}$ M and [NH$_3$] = $1.5\times10^{-2}$M. Equilibrium constant for the reaction is \underline{\hspace{2em}}.
\begin{qoptions}
\opt{A}{289}
\opt{B}{417}
\opt{C}{328}
\opt{D}{529}
\end{qoptions}

\qn{5} The molarity of 1L orthophosphoric acid (H$_3$PO$_4$) having 70\% purity by weight (specific gravity 1.54 g cm$^{-3}$) is \underline{\hspace{2em}}M.\\
(Molar mass of H$_3$PO$_4$ = 98 g mol$^{-1}$)
\begin{qoptions}
\opt{A}{55}
\opt{B}{33}
\opt{C}{22}
\opt{D}{11}
\end{qoptions}

\qn{6} \qfig{p1_chem_s1_q06_1}
An ideal gas undergoes a cyclic transformation starting from the point A and coming back to the same point by tracing the path A $\longrightarrow$ B $\longrightarrow$ C $\longrightarrow$ A as shown in the diagram. The total work done in the process is \underline{\hspace{2em}} J.
\begin{qoptions}
\opt{A}{100}
\opt{B}{200}
\opt{C}{400}
\opt{D}{50}
\end{qoptions}

\qn{7} Given below are two statements :\\
\textbf{Statement (I) :} A $\pi$ bonding MO has lower electron density above and below the inter-nuclear axis.\\
\textbf{Statement (II) :} The $\pi$* antibonding MO has a node between the nuclei.\\
In the light of the above statements, choose the most appropriate answer from the options given below:
\begin{qoptions}
\opt{A}{Both Statement I and Statement II are false}
\opt{B}{Both Statement I and Statement II are true}
\opt{C}{Statement I is false but Statement II is true}
\opt{D}{Statement I is true but Statement II is false}
\end{qoptions}

\qn{8} Which of the following compounds show colour due to d-d transition?
\begin{qoptions}
\opt{A}{CuSO$_4$.5H$_2$O}
\opt{B}{K$_2$Cr$_2$O$_7$}
\opt{C}{K$_2$CrO$_4$}
\opt{D}{KMnO$_4$}
\end{qoptions}

\qn{9} The number of radial node/s for 3p orbital is:
\begin{qoptions}
\opt{A}{1}
\opt{B}{4}
\opt{C}{2}
\opt{D}{3}
\end{qoptions}

\qn{10} Given below are two statements :\\
\textbf{Statement (I) :} Dimethyl glyoxime forms a six-membered covalent chelate when treated with NiCl$_2$ solution in presence of NH$_4$OH.\\
\textbf{Statement (II) :} Prussian blue precipitate contains iron both in (+2) and (+3) oxidation states. In the light of the above statements, choose the most appropriate answer from the options given below:
\begin{qoptions}
\opt{A}{Statement I is false but Statement II is true}
\opt{B}{Both Statement I and Statement II are true}
\opt{C}{Both Statement I and Statement II are false}
\opt{D}{Statement I is true but Statement II is false}
\end{qoptions}

\qn{11} Given below are two statements:\\
\textbf{Statement (I):} A solution of $[\mathrm{Ni(H_2O)_6}]^{2+}$ is green in colour.\\
\textbf{Statement (II):} A solution of $[\mathrm{Ni(CN)_4}]^{2-}$ is colourless.\\
In the light of the above statements, choose the most appropriate answer from the options given below:
\begin{qoptions}
\opt{A}{Both Statement I and Statement II are incorrect}
\opt{B}{Both Statement I and Statement II are correct}
\opt{C}{Statement I is incorrect but Statement II is correct}
\opt{D}{Statement I is correct but Statement II is incorrect}
\end{qoptions}

\qn{12} In case of isoelectronic species the size of F$^{-}$, Ne and Na$^{+}$ is affected by:
\begin{qoptions}
\opt{A}{Principal quantum number ($n$)}
\opt{B}{None of the factors because their size is the same}
\opt{C}{Electron-electron interaction in the outer orbitals}
\opt{D}{Nuclear charge ($z$)}
\end{qoptions}

\qn{13} According to the wave-particle duality of matter by de-Broglie, which of the following graph plot presents most appropriate relationship between wavelength of electron ($\lambda$) and momentum of electron ($p$)?
\begin{qoptions}
\opt{A}{\qfigi{p1_chem_s1_q13_1}}
\opt{B}{\qfigi{p1_chem_s1_q13_2}}
\opt{C}{\qfigi{p1_chem_s1_q13_3}}
\opt{D}{\qfigi{p1_chem_s1_q13_4}}
\end{qoptions}

\qn{14} Which of the following is strongest Bronsted base?
\begin{qoptions}
\opt{A}{\qfigi{p1_chem_s1_q14_1}}
\opt{B}{\qfigi{p1_chem_s1_q14_2}}
\opt{C}{\qfigi{p1_chem_s1_q14_3}}
\opt{D}{\qfigi{p1_chem_s1_q14_4}}
\end{qoptions}

\qn{15} The final product A, formed in the following reaction sequence is:
\qfig{p1_chem_s1_q15_1}
\begin{qoptions}
\opt{A}{Ph$-$CH$_2-$CH$_2-$CH$_3$}
\opt{B}{\qfigi{p1_chem_s1_q15_2}}
\opt{C}{\qfigi{p1_chem_s1_q15_3}}
\opt{D}{Ph$-$CH$_2-$CH$_2-$CH$_2-$OH}
\end{qoptions}

\qn{16} The major product(P) in the following reaction is :
\qfig{p1_chem_s1_q16_1}
\begin{qoptions}
\opt{A}{\qfigi{p1_chem_s1_q16_2}}
\opt{B}{\qfigi{p1_chem_s1_q16_3}}
\opt{C}{\qfigi{p1_chem_s1_q16_4}}
\opt{D}{\qfigi{p1_chem_s1_q16_5}}
\end{qoptions}

\qn{17} Identify the reagents used for the following conversion
\qfig{p1_chem_s1_q17_1}
\begin{qoptions}
\opt{A}{A = LiAlH$_4$, B = NaOH$_{(\mathrm{aq})}$,\\ C =NH$_2-$NH$_2$/KOH, ethylene glycol}
\opt{B}{A = LiAlH$_4$, B = NaOH$_{(\mathrm{alc})}$, C =Zn/HCl}
\opt{C}{A = DIBAL-H, B= NaOH$_{(\mathrm{aq})}$,\\ C = NH$_2-$NH$_2$/KOH, ethylene glycol}
\opt{D}{A = DIBAL-H, B = NaOH$_{(\mathrm{alc})}$, C = Zn/HCl}
\end{qoptions}

\qn{18} Given below are two statements:\\
\textbf{Statement-I:} High concentration of strong nucleophilic reagent with secondary alkyl halides which do not have bulky substituents will follow S$_\mathrm{N}$2 mechanism.\\
\textbf{Statement-II:} A secondary alkyl halide when treated with a large excess of ethanol follows S$_\mathrm{N}$1 mechanism.\\
In the the light of the above statements, choose the most appropriate from the questions given below:
\begin{qoptions}
\opt{A}{Statement I is true but Statement II is false}
\opt{B}{Statement I is false but Statement II is true}
\opt{C}{Both statement I and Statement II are false}
\opt{D}{Both statement I and Statement II are true}
\end{qoptions}

\qn{19} Match List I with List II
\par\nobreak
\begin{center}
\begin{tabular}{|c|p{2.2cm}|c|p{2.2cm}|}\hline
 & \multicolumn{1}{c|}{\textbf{List-I}} & & \multicolumn{1}{c|}{\textbf{List-II}}\\\hline
A. & Glucos/ NaHCO$_3$/D & I. & Gluconic acid\\\hline
B. & Glucose/ HNO$_3$ & II. & No reaction\\\hline
C. & Glucose/HI/ D & III. & n-hexane\\\hline
D. & Glucose/ Bromine water & IV. & Saccharic acid\\\hline
\end{tabular}
\end{center}
Choose the correct answer from the options given below:
\begin{qoptions}
\opt{A}{A-IV, B-I, C-III, D-II}
\opt{B}{A-II, B-IV, C-III, D-I}
\opt{C}{A-III, B-II, C-I, D-IV}
\opt{D}{A-I, B-IV, C-III, D-II}
\end{qoptions}

\qn{20} Identify A and B in the following reaction sequence.
\qfig{p1_chem_s1_q20_1}
\begin{qoptions}
\opt{A}{\qfigi{p1_chem_s1_q20_2}}
\opt{B}{\qfigi{p1_chem_s1_q20_3}}
\opt{C}{\qfigi{p1_chem_s1_q20_4}}
\opt{D}{\qfigi{p1_chem_s1_q20_5}}
\end{qoptions}

\sect{Section-II}

\qn{1} The hydrogen electrode is dipped in a solution of pH = 3 at 25$^\circ$C. The potential of the electrode will be $-$\underline{\hspace{2em}} $\times10^{-2}$ V. $\left(\dfrac{2.303\,RT}{F}=0.059\ \mathrm{V}\right)$

\qn{2} If three moles of an ideal gas at 300 K expand isothermally from 30 dm$^3$ to 45 dm$^3$ against a constant opposing pressure of 80 kPa, then the amount of heat transferred is\underline{\hspace{2em}} J.

\qn{3} Molality of 0.8 M H$_2$SO$_4$ solution (density 1.06 g cm$^{-3}$) is \underline{\hspace{2em}}$\times10^{-3}$ m.

\qn{4} Among the following oxide of p - block elements, number of oxides having amphoteric nature is\\
Cl$_2$O$_7$, CO, PbO$_2$, N$_2$O, NO, Al$_2$O$_3$, SiO$_2$, N$_2$O$_5$, SnO$_2$

\qn{5} The number of molecules/ion/s having trigonal bipyramidal shape is \dots\dots\dots .\\
PF$_5$, BrF$_5$, PCl$_5$, $[\mathrm{PtCl_4}]^{2-}$, BF$_3$, Fe(CO)$_5$

\qn{6} The `spin only' magnetic moment for $[\mathrm{Ni(NH_3)_6}]^{2+}$ is \dots.. $\times10^{-1}$ BM.\\
(Atomic number of Ni = 28)

\qn{7} Total number of compounds with Chiral carbon atoms from following is \underline{\hspace{2em}}.
\qfig{p1_chem_s2_q07_1}

\qn{8} The compound formed by the reaction of ethanal with semicarbazide contains \underline{\hspace{2em}} number of nitrogen atoms.

\qn{9} Total number of deactivating groups in aromatic electrophilic substitution reaction among the following is
\qfig{p1_chem_s2_q09_1}

\qn{10} Number of compounds which give reaction with Hinsberg's reagent is\underline{\hspace{2em}}.
\qfig{p1_chem_s2_q10_1}

\partsec{3}{MATHEMATICS}
\sect{Section-I}

\qn{1} Let 3, a, b, c be in A.P. and 3, a $-$ 1, b + 1, c + 9 be in G.P. Then, the arithmetic mean of a, b and c is :
\begin{qoptions}
\opt{A}{$-4$}
\opt{B}{$-1$}
\opt{C}{13}
\opt{D}{11}
\end{qoptions}

\qn{2} Let $\alpha$ and $\beta$ be the roots of the equation $px^2+qx-r=0$, where $p\neq0$. If $p$, $q$ and $r$ be the consecutive terms of a non-constant G.P and $\dfrac1\alpha+\dfrac1\beta=\dfrac34$, then the value of $(\alpha-\beta)^2$ is :
\begin{qoptions}
\opt{A}{$\dfrac{80}{9}$}
\opt{B}{9}
\opt{C}{$\dfrac{20}{3}$}
\opt{D}{8}
\end{qoptions}

\qn{3} If in a G.P. of 64 terms, the sum of all the terms is 7 times the sum of the odd terms of the G.P, then the common ratio of the G.P. is equal to
\begin{qoptions}
\opt{A}{7}
\opt{B}{4}
\opt{C}{5}
\opt{D}{6}
\end{qoptions}

\qn{4} A fair die is thrown until 2 appears. Then the probability, that 2 appears in even number of throws, is :
\begin{qoptions}
\opt{A}{$\dfrac56$}
\opt{B}{$\dfrac16$}
\opt{C}{$\dfrac5{11}$}
\opt{D}{$\dfrac6{11}$}
\end{qoptions}

\qn{5} If $\dfrac{{}^{11}C_1}{2}+\dfrac{{}^{11}C_2}{3}+\dots\dots+\dfrac{{}^{11}C_9}{10}=\dfrac nm$ with gcd(n, m) = 1, then n + m is equal to
\begin{qoptions}
\opt{A}{2041}
\opt{B}{2024}
\opt{C}{2025}
\opt{D}{2023}
\end{qoptions}

\qn{6} All the letters of the word ``GTWENTY'' are written in all possible ways with or without meaning and these words are written as in a dictionary. The serial number of the word ``GTWENTY'' IS
\begin{qoptions}
\opt{A}{552}
\opt{B}{551}
\opt{C}{553}
\opt{D}{1033}
\end{qoptions}

\qn{7} Let $x=x(t)$ and $y=y(t)$ be solutions of the differential equations $\dfrac{dx}{dt}+ax=0$ and $\dfrac{dy}{dt}+by=0$ respectively, $a,b\in R$. Given that $x(0)=2$; $y(0)=1$ and $3y(1)=2x(1)$, the value of $t$, for which $x(t)=y(t)$, is :
\begin{qoptions}
\opt{A}{$\log_{\frac23}2$}
\opt{B}{$\log_43$}
\opt{C}{$\log_34$}
\opt{D}{$\log_{\frac43}2$}
\end{qoptions}

\qn{8} If $(a,b)$ be the orthocentre of the triangle whose vertices are (1, 2), (2, 3) and (3, 1), and $I_1=\displaystyle\int_a^b x\sin\left(4x-x^2\right)dx$, $I_2=\displaystyle\int_a^b\sin\left(4x-x^2\right)dx$, then $36\dfrac{I_1}{I_2}$ is equal to :
\begin{qoptions}
\opt{A}{72}
\opt{B}{88}
\opt{C}{80}
\opt{D}{66}
\end{qoptions}

\qn{9} If $\sin\left(\dfrac yx\right)=\log_e|x|+\dfrac\alpha2$ is the solution of the differential equation $x\cos\left(\dfrac yx\right)\dfrac{dy}{dx}=y\cos\left(\dfrac yx\right)+x$ and $y(1)=\dfrac\pi3$, then $\alpha^2$ is equal to
\begin{qoptions}
\opt{A}{3}
\opt{B}{12}
\opt{C}{4}
\opt{D}{9}
\end{qoptions}

\qn{10} Let $a_1,a_2,\dots\dots,a_{10}$ be 10 observations such that $\displaystyle\sum_{k=1}^{10}a_k=50$ and $\displaystyle\sum_{\forall k<j}a_k\cdot a_j=1100$. Then the standard deviation of $a_1,a_2,\dots,a_{10}$ is equal to :
\begin{qoptions}
\opt{A}{5}
\opt{B}{$\sqrt5$}
\opt{C}{10}
\opt{D}{$\sqrt{115}$}
\end{qoptions}

\qn{11} The number of solutions of the equation $4\sin^2x-4\cos^3x+9-4\cos x=0$; $x\in[-2\pi,2\pi]$ is:
\begin{qoptions}
\opt{A}{1}
\opt{B}{3}
\opt{C}{2}
\opt{D}{0}
\end{qoptions}

\qn{12} Consider the relations $R_1$ and $R_2$ defined as $aR_1b\Leftrightarrow a^2+b^2=1$ for all $a,b\in R$ and $(a,b)R_2(c,d)\Leftrightarrow a+d=b+c$ for all $(a,b),(c,d)\in N\times N$. Then
\begin{qoptions}
\opt{A}{Only $R_1$ is an equivalence relation}
\opt{B}{Only $R_2$ is an equivalence relation}
\opt{C}{$R_1$ and $R_2$ both are equivalence relations}
\opt{D}{Neither $R_1$ nor $R_2$ is an equivalence relation}
\end{qoptions}

\qn{13} Consider the function $f:(0,2)\to R$ defined by $f(x)=\dfrac x2+\dfrac2x$ and the function $g(x)$ defined by\\
$g(x)=\begin{cases}\min\{f(t)\}, & 0<t\leqslant x\text{ and }0<x\leqslant1\\ \dfrac32+x, & 1<x<2\end{cases}$. Then
\begin{qoptions}
\opt{A}{$g$ is continuous but not differentiable at $x=1$}
\opt{B}{$g$ is not continuous for all $x\in(0,2)$}
\opt{C}{$g$ is neither continuous nor differentiable at $x=1$}
\opt{D}{$g$ is continuous and differentiable for all $x\in(0,2)$}
\end{qoptions}

\qn{14} Let $f:R-\left\{\dfrac{-1}{2}\right\}\to R$ and $g:R-\left\{\dfrac{-5}{2}\right\}\to R$ be defined as $f(x)=\dfrac{2x+3}{2x+1}$ and $g(x)=\dfrac{|x|+1}{2x+5}$. Then the domain of the function fog is :
\begin{qoptions}
\opt{A}{$R-\left\{-\dfrac52\right\}$}
\opt{B}{$R$}
\opt{C}{$R-\left\{-\dfrac74\right\}$}
\opt{D}{$R-\left\{-\dfrac52,-\dfrac74\right\}$}
\end{qoptions}

\qn{15} Let $g(x)=3f\left(\dfrac x3\right)+f(3-x)$ and $f''(x)>0$ for all $x\in(0,3)$. If $g$ is decreasing in $(0,\alpha)$ and increasing in $(\alpha,3)$, then $8\alpha$ is :
\begin{qoptions}
\opt{A}{24}
\opt{B}{0}
\opt{C}{18}
\opt{D}{20}
\end{qoptions}

\qn{16} If $\displaystyle\lim_{x\to0}\dfrac{3+\alpha\sin x+\beta\cos x+\log_e(1-x)}{3\tan^2x}=\dfrac13$, then $2\alpha-\beta$ is equal to :
\begin{qoptions}
\opt{A}{2}
\opt{B}{7}
\opt{C}{5}
\opt{D}{1}
\end{qoptions}

\qn{17} Let R be the interior region between the lines $3x-y+1=0$ and $x+2y-5=0$ containing the origin. The set of all values of $a$, for which the points $(a^2,a+1)$ lie in R, is :
\begin{qoptions}
\opt{A}{$(-3,-1)\cup\left(-\dfrac13,1\right)$}
\opt{B}{$(-3,0)\cup\left(\dfrac13,1\right)$}
\opt{C}{$(-3,0)\cup\left(\dfrac23,1\right)$}
\opt{D}{$(-3,-1)\cup\left(\dfrac13,1\right)$}
\end{qoptions}

\qn{18} If the foci of a hyperbola are same as that of the ellipse $\dfrac{x^2}9+\dfrac{y^2}{25}=1$ and the eccentricity of the hyperbola is $\dfrac{15}8$ times the eccentricity of the ellipse, then the smaller focal distance of the point $\left(\sqrt2,\dfrac{14}3\sqrt{\dfrac25}\right)$ on the hyperbola, is equal to
\begin{qoptions}
\opt{A}{$7\sqrt{\dfrac25}-\dfrac83$}
\opt{B}{$14\sqrt{\dfrac25}-\dfrac43$}
\opt{C}{$14\sqrt{\dfrac25}-\dfrac{16}3$}
\opt{D}{$7\sqrt{\dfrac25}+\dfrac83$}
\end{qoptions}

\qn{19} The distance of the point Q(0, 2, $-2$) form the line passing through the point P(5, $-4$, 3) and perpendicular to the lines\\
$\vec r=\left(-3\hat i+2\hat k\right)+\lambda\left(2\hat i+3\hat j+5\hat k\right)$, $\lambda\in\mathbb R$\\
and $\vec r=\left(\hat i-2\hat j+\hat k\right)+\mu\left(-\hat i+3\hat j+2\hat k\right)$, $\mu\in\mathbb R$ is :
\begin{qoptions}
\opt{A}{$\sqrt{86}$}
\opt{B}{$\sqrt{20}$}
\opt{C}{$\sqrt{54}$}
\opt{D}{$\sqrt{74}$}
\end{qoptions}

\qn{20} Let the locus of the mid points of the chords of circle $x^2+(y-1)^2=1$ drawn from the origin intersect the line $x+y=1$ at P and Q. Then, the length of PQ is :
\begin{qoptions}
\opt{A}{$\dfrac1{\sqrt2}$}
\opt{B}{$\sqrt2$}
\opt{C}{$\dfrac12$}
\opt{D}{1}
\end{qoptions}

\sect{Section-II}

\qn{1} Let the complex numbers $\alpha$ and $\dfrac1\alpha$ lie on the circles $|z-z_0|^2=4$ and $|z-z_0|^2=16$ respectively, where $z_0=1+i$. Then, the value of $100|\alpha|^2$ is\dots

\qn{2} Let A be a $2\times2$ real matrix and I be the identity matrix of order 2. If the roots of the equation $|A-xI|=0$ be $-1$ and 3, then the sum of the diagonal elements of the matrix A$^2$ is\dots\dots

\qn{3} Let the area of the region $\{(x,y):x-2y+4\geq0,\ x+2y^2\geq0,\ x+4y^2\leq8,\ y\geq0\}$ be $\dfrac mn$, where $m$ and $n$ are coprime numbers. Then $m+n$ is equal to \underline{\hspace{2em}}.

\qn{4} If the solution of the differential equation $(2x+3y-2)dx+(4x+6y-7)dy=0$, $y(0)=3$, is $ax+\beta y+3\log_e|2x+3y-\gamma|=6$, then $\alpha+2\beta+3\gamma$ is equal to \underline{\hspace{2em}}.

\qn{5} Let the area of the region $\{(x,y):0\leq x\leq3,\ 0\leq y\leq\min\{x^2+2,\,2x+2\}\}$ be A. Then 12A is equal to \underline{\hspace{2em}}.

\qn{6} Let the set of all $a\in R$ such that the equation $\cos2x+a\sin x=2a-7$ has a solution be $[p,q]$ and $r=\tan9^\mathrm{o}-\tan27^\mathrm{o}-\dfrac1{\cot63^\mathrm{o}}+\tan81^\mathrm{o}$, then $pqr$ is equal to \underline{\hspace{2em}}.

\qn{7} If $y=\dfrac{\left(\sqrt x+1\right)\left(x^2-\sqrt x\right)}{x\sqrt x+x+\sqrt x}+\dfrac1{15}\left(3\cos^2x-5\right)\cos^3x$ then $96y'\left(\dfrac\pi6\right)$ is equal to :

\qn{8} Let $S=(-1,\infty)$ and $f:S\to R$ be defined as\\
$f(x)=\displaystyle\int_{-1}^x\left(e^t-1\right)^{11}(2t-1)^5(t-2)^7(t-3)^{12}(2t-10)^{61}dt$\\
Let $p=$ Sum of square of the values of $x$, where $f(x)$ attains local maxima on S. and $q=$ Sum of the values of $x$, where $f(x)$ attains local minima on S. Then the value of $p^2+2q$ is \underline{\hspace{2em}}

\qn{9} The position vectors of the vertices A, B and C of a triangle are $2\hat i-3\hat j+3\hat k$, $2\hat i+2\hat j+3\hat k$ and $-\hat i+\hat j+3\hat k$ respectively. Let $\ell$ denotes the length of the angle bisector AD of $\angle$BAC where D is on the line segment BC, then $2\ell^2$ equals :

\qn{10} Let the line L : $\sqrt2x+y=\alpha$ pass through the point of the intersection P (in the first quadrant) of the circle $x^2+y^2=3$ and the parabola $x^2=2y$. Let the line L touch two circles C$_1$ and C$_2$ of equal radius $2\sqrt3$. If the centres Q$_1$ and Q$_2$ of the circles C$_1$ and C$_2$ lie on the y-axis, then the square of the area of the triangle PQ$_1$Q$_2$ is equal to \underline{\hspace{2em}}.

\end{multicols}
\end{document}
